% !PRINTED 379
\weberpage{379}
%<<CONT
in which the coefficients $a$ are to be regarded as known. Then from (14) we obtain four relations, if in place of $k$ we let four different constants $k_4$, $k_5$, $k_6$, $k_7$ enter:
\weq{18}{
\begin{aligned}
k_4\,(a_{4,1}\,x_1+a_{4,2}\,x_2+a_{4,3}\,x_3) &= \frac{\xi_1}{a_{4,1}}+\frac{\xi_2}{a_{4,2}}+\frac{\xi_3}{a_{4,3}} \\
k_5\,(a_{5,1}\,x_1+a_{5,2}\,x_2+a_{5,3}\,x_3) &= \frac{\xi_1}{a_{5,1}}+\frac{\xi_2}{a_{5,2}}+\frac{\xi_3}{a_{5,3}} \\
k_6\,(a_{6,1}\,x_1+a_{6,2}\,x_2+a_{6,3}\,x_3) &= \frac{\xi_1}{a_{6,1}}+\frac{\xi_2}{a_{6,2}}+\frac{\xi_3}{a_{6,3}} \\
k_7\,(a_{7,1}\,x_1+a_{7,2}\,x_2+a_{7,3}\,x_3) &= \frac{\xi_1}{a_{7,1}}+\frac{\xi_2}{a_{7,2}}+\frac{\xi_3}{a_{7,3}},
\end{aligned}
}
and now the constants $k$ are to be determined so that of these four equations one follows from the other three. The conditions for this can be formed in a symmetric manner by introducing four factors $l_4$, $l_5$, $l_6$, $l_7$, whose ratios are determined from the equations:
\weq{19}{\frac{l_4}{a_{4,i}} + \frac{l_5}{a_{5,i}} + \frac{l_6}{a_{6,i}} + \frac{l_7}{a_{7,i}} = 0, \quad i = 1, 2, 3,}
and which must then also satisfy the three equations
\weq{20}{k_4\,l_4\,a_{4,i} + k_5\,l_5\,a_{5,i} + k_6\,l_6\,a_{6,i} + k_7\,l_7\,a_{7,i} = 0, \quad i = 1, 2, 3}
from which the ratios of the $k$ are determined. A common factor of the four quantities $k$ remains, by the nature of the matter, undetermined and can be assumed arbitrarily. According to this, the functions $\xi_1$, $\xi_2$, $\xi_3$ can be rationally determined from the equations (18), and by (16) the functions $\xi_{4i}$, $\xi_{5i}$, $\xi_{6i}$, $\xi_{7i}$ are then also determined.

There remain six double tangents still wanting, which can be obtained by suitable permutations among the functions of the system (1). With this the Theorem I. placed at the head is proved.

\par\noindent\rule{0.3\textwidth}{0.4pt}\par

To see the correctness of Theorem II. as well, one need only follow our analysis backwards, regarding the coefficients $a_{k,i}$ as independent variables. Then by equations (18), (19), (20) the functions $\xi_1$, $\xi_2$, $\xi_3$ are
%CONT>>
